\section{Candidate Selection via E-Graph Extraction}\label{sec:beam}


After generating candidate abstractions,
  the \T{LLMT} algorithm invokes \T{select_library} to pick the subset of candidate patterns
  that can best be used to compress the input corpus.
%
This section describes our approach to selecting the optimal library dubbed
  \emph{targeted common subexpression elimination}.

%In this section we describe \emph{beam extraction},
%  our e-graph based approach to selecting a
%  library from the set of candidate abstractions.

%\subsection{Library Selection as E-Graph Extraction}
\mypara{Library Selection as E-Graph Extraction}
%
% At this stage (lines 6--12 in~\autoref{fig:algo-llmt}),
%   LLMT has initialized the e-graph $\mathcal{G}$
%     with the original input corpus,
%   run equality saturation on $\mathcal{G}$
%     using the given equational theory,
%   and has identified candidate abstractions
%     as a set of patterns $P$ via
%     e-graph anti-unification.
% For each candidate pattern $p \in P$
%   with parameters $\many{X}$,
%   we introduce the rewrite rule
%   $\kappa(p) \triangleq p \Rightarrow (\lambda \many{X} \bnd p)\ \many{X}$
%   as defined in \autoref{sec:au}.
%   % introduce a new rewrite rule
%   % $p \Rightarrow \T{let}\ f_p =
%   %   \lambda \many{X} \bnd p\ \T{in}\ f_p\ \many{X}$.
% LLMT uses these rewrite rules for
%   equality saturation on $\mathcal{G}$ to explore
%   all the ways the corpus may be (re-)expressed
%   using different combinations of
%   the candidate library functions (L13).
% We now desribe the next step in the algorithm which selects
%   the optimal abstractions and uses them to
%   rewrite the original corpus (L15 - L20).
Recall that candidate selection starts with an e-graph $\mathcal{G}'$,
  which represents all the ways of compressing the initial corpus and its equivalent terms
  using the candidate patterns $\patset$.
%  
We will refer to a subset $\libset \subset \patset$ as a \emph{library}.
%
The optimal size of an e-class $c$ compressed with $\libset$
  can be computed as the sum of the sizes of
  (1) the smallest term $t\in \denot{c}$
      using only the library functions in $\libset$ and
      where $\lambda$-abstractions 
      do not count toward the size of $t$, and
  (2) the smallest version of each $p \in \libset$.
      % (recall that $\lambda$-bodies themselves can be compressed using other patterns).
Note that
  the cost of defining an abstraction
  in $\libset$ is only counted once, and
  that abstractions in $\libset$ can be used
  to compress other abstractions in $\libset$.
%
Our goal is to find $\libset$ such that the root e-class compressed with $\libset$
  has the smallest size.

Given a \emph{particular} library,
  we can find the size of the smallest term
  via a relatively straightforward top-down traversal of the e-graph.
%
Hence, a na\"ive approach to library selection
  would be to enumerate all subsets of $\patset$
  and pick the one that produces the smallest term at the root. 
%
Unfortunately, this approach becomes intractable as the size of $\patset$ grows.
  % $\T{minSize}(c)$
%   represented by each e-class $c$
%   with a fixpoint approach:
%   if e-node $s() \in c$
%     (i.e., $s$ is a leaf with no children),
%     then $\T{minSize}(c) = 1$;
%   otherwise,
%     $\T{minSize}(c) = \T{min}(
%       \{1 + m_1 + \ldots + m_k |\ s(c_1, \ldots, c_k) \in c,\ m_i = \T{minSize}(c_i)\})$.
% %  if we let $T$ be
% %    $\{s(t_1, \ldots, t_k)\ |\ s(c_1, \ldots, c_k) \in c,\ t_i \in \T{minSize}(c_i)\}$
% %    then $\T{minSize}(c)$ is
% %    the subset of $T$ with minimal size.
% Because size never decreases,
%   this process terminates.
% Finding the minimal library is thus naively decidable
%   with a top-down approach:
%   simply enumerate all subsets of $\patset$,
%   compute their library cost, and
%   take the smallest.
% However, this top-down approach becomes intractable
%   once $\patset$ is moderately sized as it will have
%   exponentially many subsets.

\begin{figure}
  \textbf{E-node cost set} $\cs_N(n)$
  $$
  \begin{array}{rl}
    \cs_N(s()) &= \{ (\emptyset, 1) \} \\
    \cs_N(s(\many{a_i})) &= \{ (\libset, u + 1) \mid (\libset, u) \in \mathsf{cross}(\many{a_i}) \} \\
    % \cs_N(\T{let}\ s = \lambda \many{X} \bnd a\ \T{in}\ b) &= \mathsf{addlib}(s, \cs(a), \cs(b)) \\
    \cs_N((\lambda \many{X} \bnd a)\ b) &= \mathsf{addlib}(a, \cs(a), \cs(b)) \\
    \end{array}
  $$

  \textbf{E-class cost set} $\cs(a)$
  $$
  \begin{array}{rl}
    \cs(\{ \many{n_j} \}) &= \sfprune(\sfunify(\bigcup \many{\cs_N(n_j)})) \\
  \end{array}
  $$

  \textbf{Auxiliary definitions} $\mathsf{cross}$, $\mathsf{addlib}$, $\sfunify$, $\sfprune$
  $$
  \begin{array}{rl}
    \mathsf{cross}(z_1, z_2) &= \sfprune(\sfunify(\{ (\libset_1 \cup \libset_2, u_1 + u_2) \mid (\libset_1, u_1) \in z_1, (\libset_2, u_2) \in z_2 \})) \\
    \mathsf{addlib}(a, z_1, z_2) &= \sfprune(\sfunify(\{ (\libset_1 \cup \libset_2 \cup \{ a \}, u_2) \mid (\libset_1, u_1) \in z_1, (\libset_2, u_2) \in z_2 \})) \\
    \sfunify(z) &= \{ (\libset_1, u_1) \in z \mid \forall (\libset_2, u_2) \in z \setminus (\libset_1, u_1).\ \libset_1 \subset \libset_2 \lor u_1 < u_2 \} \\
    \sfprune(N, K, z) &=  \sftopk(\{(\libset, u) \in z \mid |\libset| \leq N \}, K)\\
  \end{array}
  $$
\caption{Cost set propagation, defined as two mutually
  recursive functions. Blue text corresponds to the partial order
  reduction optimization and red text corresponds to the beam approximation.
  $\sftopk(S, K)$ is a helper function that returns top $K$
  elements from the sorted set $S$.}
  \label{fig:costset}
\end{figure}



\mypara{Exploiting Partial Shared Structure}
%
Instead, \tool selects the optimal library using a bottom-up dynamic programming algorithm.
%
To this end, it associates each e-node and e-class with a \emph{cost set},
which is a set of pairs $(\libset, u)$
  where $\libset$ is a library and $u$ is the \emph{use cost} of this library,
  \ie the size of the smallest term represented by the e-node / e-class if it is allowed to use $\libset$
  (excluding the size of $\libset$ itself).
%
Cost sets are propagated up the e-graph using the rules shown in \autoref{fig:costset}.
%
The base case is a nullary e-node $s()$,
  which cannot use any library functions and whose size is always $1$.
%
For an e-node that has children,
\babble computes the $\mathsf{cross}$ product over the cost sets of all its child \eclasses.
%
Finally, for an application e-node $(\lambda \many{X} \bnd a)\ b$,
  the cost set \emph{must include} the library function $a$
  (in addition to some combination of libraries from the cost sets of $a$ and $b$);
  note that the use cost of the abstraction node only includes the use cost of $b$,
  since abstraction bodies are excluded from the use cost.
% Here, in addition to the cross-product between the child cost sets,
%   we need to consider the option of adding the current $\lambda$-abstraction to the children's libraries.

To compute the cost set of an e-class,
  \babble takes the union of the cost sets of all its e-nodes.
%
However, doing this naively would result in the size of the cost set growing exponentially.
%
To mitigate this, we define a \emph{partial order reduction} ($\sfunify$),
which only prunes provably sub-optimal cost sets.
%
Given two pairs $(\libset_1, u_1)$ and $(\libset_2, u_2)$ in the cost set of an e-class, 
if $\libset_1 \subset \libset_2$ and $u_1 \leq u_2$, 
then $\libset_2$ is subsumed by $\libset_1$ and can be removed from the cost set,
intuitively because $\libset_1$ can compress this e-class even better and with fewer library functions.
%
In practice this optimization prunes libraries with redundant abstraction,
where two different abstractions can be used to compress the same subterms.
% In practice, this optimization removes all abstraction
% subsets in which an abstraction isn't applied in every
% possible subterm.


% To efficiently identify a minimal-cost library, \tool determines
%   the final learned library \emph{bottom-up} during extraction,
%   rather than enumerating candidate libraries and evaluating
%   each of their costs.
% %
% Compression rewrites always introduce
%   an application of an $\lambda$-abstraction,
%   so we can derive the set of libraries used in
%   a term solely from $\lambda$ expressions it uses,
%   and thus we do not need to consider all subsets of $\patset$
%   when recursively extracting from each e-class.
% %
% Each e-class cost set contains every selection of
% libraries used by its terms, and the cost of each
% minimal term in that e-class assuming those libraries.

% \autoref{fig:costset} shows our algorithm for propagating
% the \emph{cost set} of an e-node or e-class, which represents
% the subsets of abstractions that can be
% used at the \enode or \eclass
% together with the corresponding term size in the presence of each
% subset.
% %
% The \emph{cost set} function $\cs(a)$ works similarly to $\T{minSize}$,
% with the difference being that $\cs(a)$ computes used subsets of
% abstractions on-the-fly for each smallest term, rather than setting
% it top-down.
% As shown in \autoref{fig:costset},
%   to compute the cost set
%   for an \enode that has children,
%   \babble computes the $\mathsf{cross}$ product over the libraries of all
%   children \eclasses.
% This returns a cost set by unioning the libraries
%   and adding the costs.
% There is a special case for \enodes that represent
%   an abstraction ($(\lambda \many{X} \bnd a) b$).
% To compute the cost set for these \enodes,
%   we similarly compute a cross product between the cost sets of the
%   abstraction definition and body.
% However, in this case, we also
%   add this abstraction itself to each library in the resulting cost set
%   and use the term costs of the body as the representative term costs in this \enode's cost set.

% let f = def in body
% 
% libraries(def) = {(l1, c1)}
% libraries(body) = {(l2, c2), (l4, c4)}
% 
% libraries_N(let f = def in body) = {(union(l1, l2, (f, c1)), c2), (union(l1, l4, (f, c1)), c4)}

% TODO: we could delete the para starting with \autoref{fig:costset}
% and uncomment this one instead
% Instead we can propagate information about which abstractions
% an e-node uses and its corresponding term size in the presence
% of these abstractions bottom-up, as detailed in \autoref{fig:costset}.
%
% Below we describe how \babble uses  $\cs(a)$ to
% extract the optimal set of abstractions and
% re-expresses the original term corpus in terms of this optimal set.

% First, the \verb|select_library|
% procedure from \autoref{fig:algo-llmt} finds the library selection
% with the smallest overall cost (minimizing combined library and term size)
% from the cost set of the root e-class.
% %
% \verb|select_library|
% returns a subset of the abstraction introducing rewrite rules
% %returned by anti-unification
% corresponding to the optimal library selection.
% %
% \babble uses these rewrite rules to perform
%   equality saturation (L17, \autoref{fig:algo-llmt}) which
%   introduces only the optimal set of abstractions to $\mathcal{G}$.
% Then, \verb|extract| extracts a term from the root e-class of $\mathcal{G}''$
%  which uses these abstractions,
%  breaking cycles between e-classes if needed.
% %\todo{incredibly handwavy lmao}
% This bottom-up propagation is still exponential,
% since it must track every abstraction subset used in
% a term and its subterms.
% %
% To mitigate this, we define a partial order reduction on
% cost sets, listed as $\sfunify$ in \autoref{fig:costset}.
% %
% Given two library selections $a$ and $b$ in the cost set
% of an e-class, if the libraries used in $a$ are a superset
% of those in $b$, but $a$ leads to a higher term cost than
% $b$, $a$ can be removed from the cost set.
% %
% In practice, this optimization removes all abstraction
% subsets in which an abstraction isn't applied in every
% possible subterm.

\mypara{Beam Approximation}
%
Even with the partial order reduction,
  calculating the cost set for every e-node and e-class
  can blow up exponentially.
%Even with the partial order reduction, calculating the cost
%set for every e-node and e-class in an e-graph representing a large
%program can result in exponential blowup.
%
To mitigate this,
  \tool provides the option to limit
  both the \emph{size of each library} inside a cost set
  and the \emph{size of the cost set} stored for each e-class.
%
This results in a \emph{beam-search} style algorithm,
where the cost set of each e-class is $\sfprune$ed, as shown in \autoref{fig:costset}.
%
This pruning operation first filters out libraries that have more than $N$ patterns,
then ranks the rest by the total cost (\ie the sum of use cost and the size of the library),
and finally returns the top $K$ libraries from that set.
  
  % limits the number of candidate libraries
  % (\ie a \emph{beam size})
  % and the size of each library considered at each e-class,
  % resulting in a novel beam search approach
  % to efficient, sharing-aware, e-graph extraction.
%
% This \emph{pruning} operation is listed as $\sfprune$
%   in \autoref{fig:costset}.
% It first prunes out libraries that have more than $\mathsf{n}$ abstractions,
%   then ranks the rest by the combined
%   cost of the abstractions and the term they appear in, and finally
%   returns the top $\mathsf{m}$ from that set.

%In practice, \tool provides multiple parameters which limit the number
%of candidate libraries at each e-class (\ie the \emph{beam size}),
%as well as the size of each library.
% As the next section shows, this technique is effective
%   in practice.

%def cs-prune($a, M, n$):
%    sort_by(a.libs, full-cost)
%    for a-sel in $a$:
%       if len(a.libs) >= n:
%          a.remove(a-sel)
%    retain($a$, M)
%    return sort($a$)
%\begin{figure}
\begin{lstlisting}[
    basicstyle=\footnotesize\ttfamily,
    xleftmargin=2.5em,
    numbers=left]
# Extracts a set of $n$ learned abstractions.
# $E$ is an e-graph containing the set of programs with
# all possible function definitions and applications
# $r$ is the e-class id of the root node of the e-graph
# $n$ is number of libraries to learn at a time
def choose_abstraction($E, r, n$):
    # $CS$ maps e-classes IDs to CostSets
    $CS$ = $\{\}$
    # run till fixpoint
    for $C$ in $E$.e-classes:
        for $N$ in $C$.e-nodes:
            $CS[N]$ = cs-prune(cs-unify(cs-union($CS[N]$, cs-node($N$))))
    best = $min_{full-cost}(CS[r])$
    return best.libs

\end{lstlisting}
\caption{The extraction algorithm of \babble.}
\label{fig:algo-extract}
\end{figure}

\begin{figure}
  \centering
  \begin{subfigure}{\textwidth}
\begin{lstlisting}[
    basicstyle=\footnotesize\ttfamily,
    xleftmargin=2.5em,
    numbers=left]
# The cost of an expression, e, can be expressed by a LibrarySelection,
# which is the set of libraries used in e, along with their associated costs.
# e's cost is then its AST size, not including the size of the libraries used within it.
# We define an ordering between library selections:
# comparing library selection $a$ and $b$ is defined as first comparing
# $a$.cost to $b$.cost, then comparing full_cost($a$) to full_cost(b).
datatype LibrarySelection:
    cost: int
    libs: (abstraction * cost) set

# This function computes the cost of the expression a library selection corresponds to,
# where abstraction definitions are only counted once.
# $l$ is a library selection.
def full_cost($l$):
    return $l$.cost + sum([$c$ | _, $c$ in $l$.libs])

# A CostSet is a set of possible library selections that can be used at a particular e-class.
# The set must be sorted in ascending order.
datatype CostSet:
    cost_set: LibrarySelection set
\end{lstlisting}
\end{subfigure}

\begin{subfigure}{\textwidth}
\begin{lstlisting}[
    basicstyle=\footnotesize\ttfamily,
    xleftmargin=2.5em,
    numbers=left]
# Performs partial order reduction over a CostSet, $a$.
def cs-unify($a$):
    i = 0
    while i < len($a$):
        j = i + 1
        while j < len($a$):
            if $a$[i].libs is subset of $a$[j].libs:
                remove $a$[j]
            else:
                j += 1
        i += 1
    return $a$
\end{lstlisting}
\end{subfigure}

\begin{subfigure}{\textwidth}
\begin{lstlisting}[
    basicstyle=\footnotesize\ttfamily,
    xleftmargin=2.5em,
    numbers=left]
# Performs a product between two cost sets, such that paired library selections
# have their libraries combined (via union) and their expression costs added.
# This operation is performed when an e-node has multiple children e-classes;
# the component libraries from each e-class
# must be combined via union, and their expression costs must be included.
# $a$ and $b$ are two input CostSets
  def cs-product($a, b$):
    result = $\{\}$
    for a-sel in $a$:
        for b-sel in $b$:
            # the union preserves lowest cost of abstraction
            new-libs = a-sel.libs $\cup$ b-sel.libs
            new-cost = a-sel.cost + b-sel.cost
            results.insert(new-libs, new-cost)
    return cs-unify(result)

\end{lstlisting}
\end{subfigure}
\caption{Helper functions for extraction in \babble.}
\label{fig:helper1}
\end{figure}

\begin{figure}
\begin{lstlisting}[
    basicstyle=\footnotesize\ttfamily,
    xleftmargin=2.5em,
    numbers=left]
# A variant of cost set products used when encountering an e-node that defines
# an abstraction. In this case, a product must be performed which adds an
# abstraction to all library selections in the resulting cost set.
# a is the CostSet of the abstraction definition e-class.
# e is an identifier for the abstraction.
# b is the CostSet of the rest of the expression.
def cs-add-lib($a, e, b$):
    result = $\{\}$
    for a-sel in $a$:
        for b-sel in $b$:
            # the union preserves lowest cost of abstraction
            a-size = min([ls-full-cost(l) | l in a-sel])
            new-libs = union(a-sel.libs, b-sel.libs, [($e$, a-size)])
            new-cost = b-sel.cost
            result.insert(new-libs, new-cost)
    # returns the modified cost set product of a and b
    return cs-unify(result)

# Performs a union between two CostSets while maintaining the sorted invariant.
# $a$ and $b$ are CostSet
def cs-union($a, b$):
    result = $a$
    ix = 0
    for b-elem in $b$:
        while ix < len($a$) and elem >= $a$[ix]:
            ix += 1
        result.insert(b-elem)
        ix += 1
    return result
\end{lstlisting}
\caption{Helper functions for extraction in \babble.}
\label{fig:helper2}
\end{figure}

\begin{figure}
\begin{lstlisting}[
    basicstyle=\footnotesize\ttfamily,
    xleftmargin=2.5em,
    numbers=left]
// TODO: Not sure I follow this.
# $s$ is the input e-node and $CS$ the current map from e-class ids to CostSet.
def cs-node($s, CS$):
    ($a_0, ..., a_i$) = children($s$)
    if len(children(s)) == 0:
        return [($\{\}$, 1)]
    # if $s$ introduces a library function
    elif form($s$).matches()
    let f = \x. abs in body

        if there exists a non-empty cost set for the e-classes for abs and body:
            return cs-add-lib(f, CS[abs], CS[body])
    else:
        if there exists a non-empty cost set for each of a_0, ..., a_i in CS:
            let s_0, ..., s_i = CS[a_0, ..., a_i]
            return [(ls, c + 1) | ls, c in cs-product(s_0, ..., s_i)]
    return $\{\}$

# Returns at most M library selections.
# Removes all library selections containing more than $n$ libraries.
# $a$ is the input CostSet
# $M$ is the maximum size of a CostSet
# $n$ is the number of libraries to learn at a time
def cs-prune($a, M, n$):
    sort_by(a.libs, full-cost)
    for a-sel in $a$:
       if len(a.libs) >= n:
          a.remove(a-sel)
    retain($a$, M)
    return sort($a$)
\end{lstlisting}
\caption{Helper functions for extraction in \babble.}
\label{fig:helper3}
\end{figure}

